\clearpage
\section{Scriptural Date and Location}
\label{sec:date-location}
{\small Each passage the Mass reads or sings from Scripture, once, in canonical
order; beneath each, its attribution, its setting and the dates the shared
chronology corpus holds for it. The Entrance Antiphon, both Communion antiphons
and the Alleluia verse adapt Scripture and are placed at the verses they adapt;
the three orations carry no biblical date.\par}

\begin{dossiertable}
Responsorial Psalm & Ps 23 (22):1--3a, 3b--4, 5, 6; response v.~6cd & No place
named in the verses &
\chronodate{responsorial-psalm}{\chronologyannotation{responsorial-psalm}}\\
\cmidrule(lr){1-4}
\dossierprose{Inherited attribution: the Vulgate title \latin{Psalmus David},
in the Douay--Rheims ``A psalm for David'' (22:1), ascribes the psalm to David;
the title gives no date or occasion, and the corpus holds no date for the
ascription. The setting answer is unresolved because the psalm names no
occasion. The composition date shown is the Psalter's common critical bound,
taken from the introduction to the Psalms in the New American Bible Revised
Edition, which holds that no individual psalm can be dated securely; it is not
a date for this psalm.}
\midrule
Communion Antiphon, first option & cf.\ Ps 34 (33):11, an adaptation & By the title,
the court of Geth, where David fled from Saul to Achis (1~Kgs 21:10) &
\chronodate{communion-antiphon-a}{\chronologyannotation{communion-antiphon-a}}\\
\cmidrule(lr){1-4}
\dossierevent{Narrated event, by the title: David changing his countenance
before Achimelech (33:1); the Douay--Rheims refers the title to 1~Kings 21,
where David feigns madness before Achis, king of Geth (21:10--15).}
\cmidrule(lr){1-4}
\dossierprose{Inherited attribution: the title ``For David'' (33:1). The
attribution answer is David's reign in the usual chronology reported by the
\work{Catholic Encyclopedia}'s article on David (Corbett), which sets it beside
later datings; the corpus marks it disputed. The superscription answer dates
the episode the title names, on the A.M. reckoning that Haydock's notes report;
it is the title's own assertion and not a date of composition. The composition
answer is the Psalter's common bound, not a date for this psalm. The antiphon's
``rich'' is the Latin and Greek reading; the Hebrew has young lions.}
\midrule
Entrance Antiphon & Ps 130 (129):3--4, an adaptation & No place named in the
verses &
\chronodate{entrance-antiphon}{\chronologyannotation{entrance-antiphon}}\\
\cmidrule(lr){1-4}
\dossierprose{A song of ascents, titled \latin{Canticum graduum}, ``A gradual
canticle''; the title names no author and no occasion, and the setting answer
is unresolved. The composition answer is the Psalter's common bound, not a date
for this psalm. The antiphon uses v.~3 and the first clause of v.~4; the query
reaches the whole of both verses.}
\midrule
First Reading & Is 25:6--10a & ``In this mountain'' (25:6, 7, 10), the
``mount Sion'' and Jerusalem where the Lord of hosts reigns (24:23); the vision
``concerning Juda and Jerusalem'' (1:1) &
\chronodate{first-reading}{\chronologyannotation{first-reading}}\\
\cmidrule(lr){1-4}
\dossierprose{Traditional attribution: ``Isaias the son of Amos'' (1:1). The
answer is the length of the prophet's ministry as the \work{Catholic
Encyclopedia}'s article on Isaias (Souvay) gives it, from the last year of
King Ozias to the latest prophecies it admits; the same article reports an old
tradition that the prophet lived into the reign of Manasses, which the answer
does not carry. It dates the prophet, not chapters 24--27 or this oracle, for
which the corpus holds no composition date. The query for v.~10 also covers
the clause on Moab, which is not read.}
\midrule
% booktabs leaves a legal page break after each partial rule, which no
% row-end hint reaches; the Gospel's two forms and their notes start a
% page together.
\pagebreak
Gospel, longer form & Mt 22:1--14 & In the temple at Jerusalem, to ``the chief
priests and ancients of the people'' (21:23) &
\chronodate{gospel-long}{\chronologyannotation{gospel-long}}\\
Gospel, shorter form & Mt 22:1--10 & As for the longer form &
\chronodate{gospel-short}{\chronologyannotation{gospel-short}}\\
\cmidrule(lr){1-4}
\dossierevent{Narrated event: a parable spoken in the temple, after the
question about Jesus' authority (21:23--27) and the parables of the two sons
and of the vineyard.}
\cmidrule(lr){1-4}
\dossierprose{Traditional attribution: St~Matthew the Apostle. The event answer
dates the telling of the parable, not the wedding it narrates: it is a
derivation within the chronology of the \work{Catholic Encyclopedia}'s article
on Jesus Christ (Maas), which places the parable of the marriage feast on the
Tuesday of the last week in Jerusalem. The composition answers keep as
disputed alternatives the figures the \work{Catholic Encyclopedia} reports. Its
article on the Gospel (Jacquier) gives four from the ancient writers, who ``are
at variance'', and the years that ``Catholic critics, in general, favour''
where ``opinion is rather divided''; the last figure is said in its article on
the New Testament (Durand) of the original text of Matthew, which that author
holds to have been Aramaic. These are dates of the book, not of the saying.}
\midrule
Alleluia verse & cf.\ Eph 1:17--18, an adaptation & From Paul, ``the prisoner
of Jesus Christ'' (3:1); to ``all the saints who are at Ephesus'' (1:1) &
\chronodate{gospel-acclamation}{\chronologyannotation{gospel-acclamation}}\\
\cmidrule(lr){1-4}
\dossierprose{Traditional attribution: St~Paul, writing in captivity. The first
range is the one the \work{Catholic Encyclopedia}'s article on the letter gives
for the letters of the captivity, where it leaves the place of the captivity,
Caesarea or Rome, ``a much mooted question''; the second figure is the year its
chronology of St~Paul (Prat) gives those letters. They date the letter, not
the liturgical verse made from it.}
\midrule
Second Reading & Phil 4:12--14, 19--20 & From Paul ``in my bands'' (1:7, 13),
with greetings from ``Caesar's household'' (4:22); to ``all the saints in
Christ Jesus who are at Philippi'' (1:1) &
\chronodate{second-reading}{\chronologyannotation{second-reading}}\\
\cmidrule(lr){1-4}
\dossierprose{Traditional attribution: St~Paul, writing in captivity. The first
figure is the year the \work{Catholic Encyclopedia}'s chronology of St~Paul
(Prat) gives the letters of the captivity, within his Roman captivity; the
second is the range its article on the letter (Vander Heeren) gives for Paul
``at Rome'', which it calls ``the nearly unanimous opinion''.}
\midrule
Communion Antiphon, second option & 1~Jn 3:2, an adaptation & No place named
in the letter &
\chronodate{communion-antiphon-b}{\chronologyannotation{communion-antiphon-b}}\\
\cmidrule(lr){1-4}
\dossierprose{Traditional attribution: St~John the Apostle. The range is the one
the \work{Catholic Encyclopedia}'s article on the New Testament (Durand) gives
for ``the Johannine writings''; its article on the Epistles of St~John (Drum)
finds the arguments ``probable in favour of Ephesus and also for the last few
years of the first century''.}
\end{dossiertable}
\clearpage
